% =====================================================================
\section{Thermodynamic interpretation}\label{sec:thermo}
% =====================================================================

We now read the framework of \cref{sec:framework} in the language of
the thermodynamics of internal variables, identifying the energy
balance with the first law and the dissipation inequality with the
second law~\cite{germain1983,halphen1975generalized}.

\subsection{Free energy and the first law}\label{sec:first_law}

The state of the system is the pair $(\bs q, S)$, and $E$ plays the role
of a \emph{Helmholtz-type free energy}: the portion of the energy
stored in the configuration and recoverable through quasi-static
changes,
\begin{equation}\label{eq:free_energy}
  E(\bs q, S) = \underbrace{F(\bs q)}_{\text{configurational}}
  \;-\;\underbrace{\bs q\cdot\bs\ell(S)}_{\text{coupling to environment}}.
\end{equation}
The configurational term $F$ is intrinsic and independent of the
environment; the coupling term $-\bs q\cdot\bs\ell(S)$ acts as a
generalised chemical potential that biases the state. The conjugate
driving force is $\bs f=-\partial_{\bs q}E$, as in~\eqref{eq:force}, and
the power injected by the loading at frozen state is
$\mathcal P_S=\partial_S E\,\dot S$, as in~\eqref{eq:power}.

Rearranging the energy balance~\eqref{eq:energy_balance} over $[0,T]$
yields the \emph{first law}
\begin{equation}\label{eq:first_law}
  \underbrace{\Delta E}_{\text{change in free energy}}
  \;+\;
  \underbrace{\Diss_\Psi(\bs q;[0,T])}_{\mathcal D_{\mathrm{tot}}\ \ge 0}
  \;=\;
  \underbrace{\int_0^T \partial_S E(\bs q,S)\,\dot S\dd t}_{\mathcal W\ \text{(work of the environment)}},
\end{equation}
with $\Delta E=E(\bs q(T),S(T))-E(\bs q(0),S(0))$: the work supplied by
the environment is stored as free energy or irreversibly dissipated.

\subsection{Entropy production and the second law}\label{sec:second_law}

The instantaneous \emph{dissipation rate} is
\begin{equation}\label{eq:diss_rate}
  \mathcal D(t) := \bs f(\bs q(t),S(t))\cdot\dot{\bs q}(t).
\end{equation}

\begin{proposition}[Second law]\label{prop:second_law}
Along any solution of~\eqref{eq:biot} one has $\mathcal D(t)\ge 0$ for
a.e.\ $t$. Moreover the \emph{excess dissipation}
$\Xi(t):=\bs f(t)\cdot\dot{\bs q}(t)-\Psi(\dot{\bs q}(t))$ vanishes
identically, $\Xi\equiv 0$.
\end{proposition}

This is an immediate consequence of the saturation
identity~\eqref{eq:saturation_solution} (see \cref{app:framework}).
The non-negativity $\mathcal D\ge 0$ is the second law: dissipation is
non-negative regardless of the direction of evolution. The identity
$\Xi\equiv 0$ expresses a \emph{principle of minimal (economical)
dissipation}: a rate-independent system saturates the dissipation
inequality, dissipating exactly the minimum compatible with the energy
balance, with no viscous excess. In an isothermal setting at absolute
temperature $\Theta$, the irreversible entropy production rate is
$\Theta\,\dot s_{\mathrm{irr}}=\mathcal D(t)\ge0$, so that over the
whole process
\begin{equation}\label{eq:entropy}
  \Delta s_{\mathrm{irr}} = \frac{\mathcal D_{\mathrm{tot}}}{\Theta}
  = \frac{\Diss_\Psi(\bs q;[0,T])}{\Theta}\ \ge\ 0,
\end{equation}
the standard form of the second law.

\subsection{Geometric dissipation over closed cycles}\label{sec:hysteresis}

A hallmark of rate-independent thermodynamics is that the energy
dissipated over a closed loading cycle depends only on the geometry of
the trajectory, not on its time scale. Consider a cycle with
$S(0)=S(T)=S_0$ that returns the state to its initial value, so that
$\Delta E=0$. Integrating~\eqref{eq:first_law} and using
$\mathcal D\ge0$ (\cref{prop:second_law}) gives
\begin{equation}\label{eq:cycle}
  \mathcal W = \oint_0^T \partial_S E(\bs q,S)\,\dot S\dd t
  = \mathcal D_{\mathrm{tot}} \ \ge\ 0.
\end{equation}
Hence the work performed by the environment over a closed cycle cannot
be recovered and is entirely converted into irreversible entropy
production. In the $(\bs q,\bs f)$ plane this manifests as a
non-vanishing \emph{hysteresis loop} whose enclosed area equals exactly
$\mathcal D_{\mathrm{tot}}$; since $\bs f=\bs\ell(S)-\bs g(\bs q)$
(\cref{eq:force}), the same loop, up to the fixed reparametrisation
$\bs f\leftrightarrow\bs\ell(S)$, is the one plotted directly in the
$(\ell,q)$ loading--state plane in the numerical experiments of
\cref{sec:experiments}. The rate-independent structure therefore
guarantees that returning the environment to its baseline leaves a
permanent thermodynamic imprint---a property that, in the motivating
biological application, models the non-volatility of the epigenetic
state.
